We use the following notation to make the methodology more concise. For
subject $i$, let
\[
    s_i(\boldsymbol\gamma)
    =
    1+\mathbf{x}_i^\top\boldsymbol\gamma,
    \qquad
    g_i(\boldsymbol\gamma)
    =
    I\{s_i(\boldsymbol\gamma)\geq 0\},
\]
and define
\[
    Z_i(\boldsymbol\gamma)
    =
    \left(
        a_i,
        g_i(\boldsymbol\gamma),
        a_i g_i(\boldsymbol\gamma)
    \right)^\top,
\]
where $a_i$ is the treatment indicator and $\mathbf{x}_i$ is the vector of
$p$ biomarkers. The Cox partial likelihood, conditional on
$\boldsymbol\gamma$, is
\begin{equation}
\label{eq:partialLikFixedThreshold}
    L(\boldsymbol\beta\mid\boldsymbol\gamma)
    =
    \prod_{i=1}^{n}
    \left[
        \frac{
            \exp\{Z_i(\boldsymbol\gamma)^\top\boldsymbol\beta\}
        }{
            \displaystyle\sum_{j\in R(t_i)}
            \exp\{Z_j(\boldsymbol\gamma)^\top\boldsymbol\beta\}
        }
    \right]^{\delta_i},
\end{equation}
where $R(t_i)$ is the index set of subjects at risk immediately before time
$t_i$, and $\delta_i$ is the event indicator.

We assign independent normal priors to the components of
$\boldsymbol\gamma$. In particular,
\[
    \gamma_\ell\sim N(\mu_\ell,\sigma_\ell^2),
    \qquad \ell=1,\ldots,p,
\]
so that, with
$\boldsymbol\mu=(\mu_1,\ldots,\mu_p)^\top$ and
$\Sigma=\operatorname{diag}(\sigma_1^2,\ldots,\sigma_p^2)$,
\begin{equation}
\label{eq:gammaPriorFixedThreshold}
    L(\boldsymbol\gamma\mid\boldsymbol\mu,\Sigma)
    =
    \prod_{\ell=1}^{p}
    \frac{1}{\sqrt{2\pi\sigma_\ell^2}}
    \exp\left\{
        -\frac{(\gamma_\ell-\mu_\ell)^2}{2\sigma_\ell^2}
    \right\}.
\end{equation}
The common-variance specification follows by setting
$\sigma_\ell^2=\sigma^2$ for every $\ell$. Under a flat prior for
$\boldsymbol\beta$, the resulting posterior based on the Cox partial
likelihood is
\begin{equation}
\label{eq:posteriorFixedThreshold}
    \pi(\boldsymbol\beta,\boldsymbol\gamma\mid\textbf{data})
    \propto
    L(\boldsymbol\beta\mid\boldsymbol\gamma)
    L(\boldsymbol\gamma\mid\boldsymbol\mu,\Sigma).
\end{equation}
When biomarker shrinkage is desired, we use the penalized target
\begin{equation}
\label{eq:penalizedPosteriorFixedThreshold}
    \pi_\lambda(\boldsymbol\beta,\boldsymbol\gamma
    \mid\textbf{data})
    \propto
    L(\boldsymbol\beta\mid\boldsymbol\gamma)
    L(\boldsymbol\gamma\mid\boldsymbol\mu,\Sigma)
    \exp\left\{-\lambda\sum_{\ell=1}^{p}|\gamma_\ell|\right\},
    \qquad \lambda\geq 0.
\end{equation}
Thus, $\lambda=0$ gives the unpenalized posterior in
Equation~\ref{eq:posteriorFixedThreshold}, whereas larger values of
$\lambda$ produce stronger shrinkage of the biomarker coefficients.

We sample from Equation~\ref{eq:penalizedPosteriorFixedThreshold} using a
Metropolis-within-Gibbs algorithm. Let $R$ denote the number of retained
posterior draws, $B$ the number of burn-in iterations, and $TN$ the thinning
interval. The sampler therefore performs $B+R\times TN$ iterations. The first
$B$ draws are discarded and, thereafter, every $TN$th draw is retained.
Algorithm~\ref{al:mcmcAlgo} outlines the sampling procedure.

\begin{algorithm}[p]
\large
\SetAlgoLined
\DontPrintSemicolon

\textbf{Initialisation:}
Draw or set $\boldsymbol\gamma_0$. Given $\boldsymbol\gamma_0$, initialise
$\boldsymbol\beta_0$ by fitting the Cox model with covariates
$Z_i(\boldsymbol\gamma_0)$\;

\For{$k=0,\ldots,B+R\times TN-1$}{
    Set $\widetilde{\boldsymbol\gamma}^{(0)}=\boldsymbol\gamma_k$\;

    \For{$\ell=1,\ldots,p$}{
        Propose
        $\gamma_\ell^*\sim
        N(\widetilde{\gamma}^{(\ell-1)}_\ell,\tau_\ell^2)$ and form
        $\boldsymbol\gamma^*$ by replacing the $\ell$th component of
        $\widetilde{\boldsymbol\gamma}^{(\ell-1)}$ with
        $\gamma_\ell^*$\;

        Define
        \[
            h(\boldsymbol\gamma\mid\boldsymbol\beta_k)
            =
            \log L(\boldsymbol\beta_k\mid\boldsymbol\gamma)
            +
            \sum_{r=1}^{p}\log
            \phi(\gamma_r;\mu_r,\sigma_r^2)
            -
            \lambda\sum_{r=1}^{p}|\gamma_r|,
        \]
        where $\phi(\,\cdot\,;\mu_r,\sigma_r^2)$ is the normal prior
        density\;

        Accept $\boldsymbol\gamma^*$ with probability
        \[
            \alpha_{\gamma,\ell}
            =
            \min\left\{
                1,
                \exp\left[
                    h(\boldsymbol\gamma^*\mid\boldsymbol\beta_k)
                    -
                    h(\widetilde{\boldsymbol\gamma}^{(\ell-1)}
                    \mid\boldsymbol\beta_k)
                \right]
            \right\}.
        \]

        If the proposal is accepted, set
        $\widetilde{\boldsymbol\gamma}^{(\ell)}
        =\boldsymbol\gamma^*$; otherwise, set
        $\widetilde{\boldsymbol\gamma}^{(\ell)}
        =\widetilde{\boldsymbol\gamma}^{(\ell-1)}$\;
    }

    Set $\boldsymbol\gamma^\dagger
    =\widetilde{\boldsymbol\gamma}^{(p)}$. Given
    $\boldsymbol\gamma^\dagger$, fit the Cox model with covariates
    $Z_i(\boldsymbol\gamma^\dagger)$ and obtain the covariance estimate
    $\widehat{\Sigma}_{k+1}$\;

    \eIf{the Cox fit converges}{
        Set $\boldsymbol\gamma_{k+1}=\boldsymbol\gamma^\dagger$\;

        \eIf{$\widehat{\Sigma}_{k+1}$ is finite}{
            Propose
            $\boldsymbol\beta^*\sim
            N(\boldsymbol\beta_k,\widehat{\Sigma}_{k+1})$\;

            Accept $\boldsymbol\beta^*$ with probability
            \[
                \alpha_\beta
                =
                \min\left\{
                    1,
                    \frac{
                        L(\boldsymbol\beta^*\mid\boldsymbol\gamma^\dagger)
                    }{
                        L(\boldsymbol\beta_k\mid\boldsymbol\gamma^\dagger)
                    }
                \right\}.
            \]

            Set $\boldsymbol\beta_{k+1}=\boldsymbol\beta^*$ if the proposal
            is accepted; otherwise set
            $\boldsymbol\beta_{k+1}=\boldsymbol\beta_k$\;
        }{
            Set $\boldsymbol\beta_{k+1}=\boldsymbol\beta_k$\;
        }
    }{
        Set $(\boldsymbol\beta_{k+1},\boldsymbol\gamma_{k+1})
        =(\boldsymbol\beta_k,\boldsymbol\gamma_k)$\;
    }

    After burn-in, retain every $TN$th draw until $R$ posterior samples are
    stored\;
}

\caption{Metropolis--Hastings sampler for the fixed-threshold single-index
Cox model.}
\KwResult{Posterior samples of $\boldsymbol\beta$,
$\boldsymbol\gamma$, and the log partial likelihood.}
\label{al:mcmcAlgo}
\end{algorithm}

We give a detailed explanation of the steps in
Algorithm~\ref{al:mcmcAlgo} below.

\textbf{Parameter initialisation:}
The initial value $\boldsymbol\gamma_0$ may be supplied using prior scientific
knowledge or drawn from the normal prior in
Equation~\ref{eq:gammaPriorFixedThreshold}. Once
$\boldsymbol\gamma_0$ has been selected, the subgroup indicators
$g_i(\boldsymbol\gamma_0)$ and design vectors
$Z_i(\boldsymbol\gamma_0)$ are computed. The initial value
$\boldsymbol\beta_0$ is then obtained from the partial likelihood estimate in
the Cox model with these design vectors. In this formulation, no separate
cut-point is updated: the threshold is fixed at zero and the intercept in the
single-index score is fixed at one.

\textbf{Step 1:}
Each component of $\boldsymbol\gamma$ is updated in turn using a normal
random-walk proposal. Here, $\tau_\ell^2$ is the proposal variance and is
distinct from the prior variance $\sigma_\ell^2$. The proposal is centered at
the current value of the $\ell$th component, whereas $\mu_\ell$ is the fixed
prior mean. Because the random-walk proposal is symmetric, its density cancels
from the Metropolis--Hastings ratio. The acceptance probability therefore
depends only on the difference between the proposed and current values of the
penalized log-posterior kernel $h$. The subgroup indicators, and hence the Cox
partial likelihood, are recomputed after each proposed component update.

\textbf{Step 2:}
After all components of $\boldsymbol\gamma$ have been updated, a Cox model is
fit using $Z_i(\boldsymbol\gamma^\dagger)$. Its estimated covariance matrix is
used to scale a multivariate normal random-walk proposal for
$\boldsymbol\beta$. Conditional on $\boldsymbol\gamma^\dagger$, this proposal
is symmetric and the prior on $\boldsymbol\beta$ is flat, so the acceptance
probability reduces to the ratio of Cox partial likelihoods shown in the
algorithm. If the Cox fit fails to converge, both parameter vectors are
returned to their values at the start of the iteration. If the Cox fit
succeeds but its covariance estimate is unusable, or if the proposed
$\boldsymbol\beta^*$ is rejected, the updated
$\boldsymbol\gamma^\dagger$ is retained and $\boldsymbol\beta$ remains
unchanged.

\subsection{Estimation}

Let
$\{(\boldsymbol\beta^{(r)},\boldsymbol\gamma^{(r)}):r=1,\ldots,R\}$
denote the retained posterior draws after burn-in and thinning. We use the
posterior means as the default point estimates,
\begin{equation}
\label{eq:posteriorMeanEstimates}
    \widehat{\boldsymbol\beta}
    =
    \frac{1}{R}\sum_{r=1}^{R}\boldsymbol\beta^{(r)},
    \qquad
    \widehat{\boldsymbol\gamma}
    =
    \frac{1}{R}\sum_{r=1}^{R}\boldsymbol\gamma^{(r)}.
\end{equation}
Posterior medians may be reported instead when a marginal posterior is notably
skewed.

Equal-tailed credible intervals are obtained directly from the empirical
posterior distribution. For any scalar parameter $\eta$, let
\[
    \widehat F_R(w)
    =
    \frac{1}{R}\sum_{r=1}^{R}I\{\eta^{(r)}\leq w\},
    \qquad
    \widehat Q_R(u)
    =
    \inf\{w:\widehat F_R(w)\geq u\}.
\]
The $100(1-\alpha)\%$ equal-tailed credible interval is then
\begin{equation}
\label{eq:credibleIntervalFixedThreshold}
    \left[
        \widehat Q_R\left(\frac{\alpha}{2}\right),
        \widehat Q_R\left(1-\frac{\alpha}{2}\right)
    \right].
\end{equation}
Equation~\ref{eq:credibleIntervalFixedThreshold} is applied componentwise to
$\boldsymbol\beta$ and $\boldsymbol\gamma$.

Because subgroup membership is a nonlinear function of
$\boldsymbol\gamma$, its uncertainty is summarized from the posterior draws
rather than only by substituting $\widehat{\boldsymbol\gamma}$ into the index.
For subject $i$, the posterior probability of membership in the
index-positive subgroup is estimated by
\begin{equation}
\label{eq:posteriorSubgroupProbability}
    \widehat P_i
    =
    \frac{1}{R}\sum_{r=1}^{R}
    I\left\{1+\mathbf{x}_i^\top\boldsymbol\gamma^{(r)}\geq 0\right\}.
\end{equation}
A subject may be assigned to this subgroup when
$\widehat P_i\geq 0.5$, while values near $0.5$ indicate substantial
classification uncertainty.

Writing $\boldsymbol\beta=(\beta_1,\beta_2,\beta_3)^\top$, the treatment
hazard ratio is $\exp(\beta_1)$ for the index-negative subgroup and
$\exp(\beta_1+\beta_3)$ for the index-positive subgroup. Point estimates and
credible intervals for these hazard ratios are obtained by applying the
corresponding transformation to every retained posterior draw and then taking
the empirical mean or median and quantiles. This draw-by-draw transformation
preserves posterior dependence between $\beta_1$ and $\beta_3$.
